% Options for packages loaded elsewhere
\PassOptionsToPackage{unicode}{hyperref}
\PassOptionsToPackage{hyphens}{url}
%
\documentclass[
  10pt,
]{article}
\usepackage{amsmath,amssymb}
\usepackage{iftex}
\ifPDFTeX
  \usepackage[T1]{fontenc}
  \usepackage[utf8]{inputenc}
  \usepackage{textcomp} % provide euro and other symbols
\else % if luatex or xetex
  \usepackage{unicode-math} % this also loads fontspec
  \defaultfontfeatures{Scale=MatchLowercase}
  \defaultfontfeatures[\rmfamily]{Ligatures=TeX,Scale=1}
\fi
\usepackage{lmodern}
\ifPDFTeX\else
  % xetex/luatex font selection
\fi
% Use upquote if available, for straight quotes in verbatim environments
\IfFileExists{upquote.sty}{\usepackage{upquote}}{}
\IfFileExists{microtype.sty}{% use microtype if available
  \usepackage[]{microtype}
  \UseMicrotypeSet[protrusion]{basicmath} % disable protrusion for tt fonts
}{}
\makeatletter
\@ifundefined{KOMAClassName}{% if non-KOMA class
  \IfFileExists{parskip.sty}{%
    \usepackage{parskip}
  }{% else
    \setlength{\parindent}{0pt}
    \setlength{\parskip}{6pt plus 2pt minus 1pt}}
}{% if KOMA class
  \KOMAoptions{parskip=half}}
\makeatother
\usepackage{xcolor}
\usepackage[margin=0.85in]{geometry}
\setlength{\emergencystretch}{3em} % prevent overfull lines
\providecommand{\tightlist}{%
  \setlength{\itemsep}{0pt}\setlength{\parskip}{0pt}}
\setcounter{secnumdepth}{-\maxdimen} % remove section numbering
\usepackage{amsmath,amssymb,amsthm}
\usepackage{longtable}
\ifLuaTeX
  \usepackage{selnolig}  % disable illegal ligatures
\fi
\usepackage{bookmark}
\IfFileExists{xurl.sty}{\usepackage{xurl}}{} % add URL line breaks if available
\urlstyle{same}
\hypersetup{
  pdftitle={Divisor-relative finite-energy equations for every big rational line bundle},
  pdfauthor={Unverified research draft for original P16.5},
  hidelinks,
  pdfcreator={LaTeX via pandoc}}

\title{Divisor-relative finite-energy equations for every big rational
line bundle}
\author{Unverified research draft for original P16.5}
\date{5 October 2026}

\begin{document}
\maketitle

\section{Abstract}\label{abstract}

We reconstruct a relative equation theory on every actual big rational
line bundle on a smooth integral projective variety over every
algebraically closed trivially valued field of characteristic zero. The
retained datum is an actual Kodaira decomposition \(L=A+D\), with \(A\)
ample and \(D\) effective rational Cartier. The reference is the actual
divisor potential \(\sigma=\log|s_D|\). We prove that its entire genuine
reference-dominated cone is \(\sigma+\operatorname{PSH}(A)\), including
all points of the divisor. Division is repaired by compactness rather
than an incorrect monotonicity assertion about divided approximants. The
inverse quotient is a uniquely determined actual psh extension on the
original variety. The relative energy is an integral of this quotient
against genuine relative metric products. We identify these products
first locally, with the correct sign in the frame identity, and then
prove their nonpluripolar extension, raw mass \(A^n\), unbounded and
weighted decreasing limits, comparison, intrinsic target criterion,
actual attainment and uniqueness modulo constants. Every continuous
direction on the original space is available with both derivative signs.
The ample equation is used solely in its exact task-granted actual ample
domain.

The additional zero-volume equation proved here is exactly the actual
trivial-bundle equation in positive dimension. We distinguish a big
blowup with nonpositive self-intersection from a genuinely zero-volume
example with noninjective top operator. This is a theorem about
prescribed retained Kodaira references, not a class-volume or
minimal-singularity theorem, and not a general zero-volume equation. All
original big rational bundles occur; no finite-generation, cardinality,
invariance or measure-support hypothesis replaces that domain. The
conclusion is a mathematical proof draft relative to an explicit task
grant and published premises. Evidence processing and compilation do not
certify it.

\section{1. Statement and the exact retained
domain}\label{statement-and-the-exact-retained-domain}

Let \(k\) be algebraically closed, of characteristic zero, with the
trivial absolute value. It is complete and may have arbitrary
cardinality. Let \(X\) be a smooth integral projective \(k\)-variety of
dimension \(n\geq1\). Analytification and measures are Berkovich
analytification and positive Radon measures. A rational line bundle
means an actual element of \(\operatorname{Pic}(X)_{\mathbb Q}\),
cleared by a common positive integer whenever frames and sections are
used. Potentials are scalar differences from the canonical trivial
metric. We retain the actual rational-bundle isomorphism \[
             L\simeq A\otimes\mathcal O_X(D),                  \tag{1.1}
\] where \(A\) is ample and \(D\) is effective rational Cartier. This
datum is fixed before any measure is specified. Put \[
 Z=|D|^{\rm an},\quad U=X^{\rm an}\setminus Z,\quad
 \sigma=m^{-1}\log\|s_{mD}\|_{\rm triv},\quad V_D=(A^n)>0.
                                                               \tag{1.2}
\] The definition of \(\sigma\) is independent of the clearing integer.
It is finite on \(U\), equals \(-\infty\) on \(Z\), is at most zero, and
equals zero at the generic trivial valuation \(v_{\rm triv}\). Every
function-field divisorial valuation belongs to \(U\) and has finite
\(\sigma\).

We use the genuine BJ cone
\(\operatorname{PSH}(X,L)=\operatorname{PSH}(c_1(L))\), not a newly
declared cone of formal symbols. It consists of generically finite usc
decreasing-net limits of generically finite Fubini--Study potentials in
rational classes tending to \(c_1(L)\). For an ample actual rational
bundle these are decreasing same-class genuine approximants, with the
fixed-bundle conversion stated in the task grant. Define the geometric
reference cone \[
 \mathcal P_D(L)=\{u\in\operatorname{PSH}(X,L):u\leq\sigma+C
                                      \text{ for some }C\in\mathbb R\}. \tag{1.3}
\] The topology on this cone is pointwise convergence on function-field
divisorial valuations. This is the usual weak psh topology; it is not
pointwise convergence at all analytic points.

\textbf{Theorem 1 (full big domain with retained divisor reference).}
Every actual big rational \(L\) on \(X\) admits (1.1). For every such
datum the map \[
 F_D:\operatorname{PSH}(X,A)\longrightarrow\mathcal P_D(L),
                 \qquad F_D(\psi)=\sigma+\psi                  \tag{1.4}
\] is an order-preserving homeomorphism, whose inverse \(Q_D\) is
subtraction on \(U\) followed by a uniquely determined actual \(A\)-psh
extension on all of \(X^{\rm an}\). Suprema and constants are preserved.
The following constructions and conclusions hold on the original
carrier.

\begin{enumerate}
\def\labelenumi{\arabic{enumi}.}
\item
  Relative bounded potentials mean \(Q_Du\) bounded. Their genuine local
  metric intersection products on \(U\) have a canonical nonpluripolar
  extension. Relative energy, initially an actual integral against these
  products and then extended by reference truncations, satisfies \[
   E_D(u)=E_{{\rm raw},A}(Q_Du),\qquad
   \mathcal E_D^1(L)=\{u\in\mathcal P_D(L):E_D(u)>-\infty\}
                      =F_D(\mathcal E^1(X,A)).                \tag{1.5}
  \] These are actual functions or metrics; no completion is their
  definition.
\item
  For \(u_i\in\mathcal E_D^1(L)\), the genuine relative mixed metric
  product is \[
   \mathcal M_D(u_1,\ldots,u_n)
      =\bigwedge_{i=1}^n(c_1(A)+dd^c Q_Du_i).                 \tag{1.6}
  \] This formula is a proved geometric identity, not the definition of
  an inverse measure. The left side is constructed before a target is
  chosen, by the local metric product and reference-relative
  nonpluripolar truncation. It is positive Radon, symmetric, has raw
  mass \(V_D\), and gives zero mass to the fixed divisor and to every
  pluripolar set. It has finite-energy locality, comparison and
  domination. It is continuous under genuine unbounded decreasing nets
  with finite-energy endpoints, including the precisely stated
  integrable quotient-weighted limits in Section 6.
\item
  The intrinsic admissible targets are precisely positive Radon \(\mu\)
  such that \[
   \mu(X^{\rm an})=V_D,\qquad
   D_D(\mu):=\sup_{u\in\mathcal E_D^1(L)}
           \left(E_D(u)-\int Q_Du\,d\mu\right)<\infty.        \tag{1.7}
  \] The energy condition is equivalent to integrability of every actual
  quotient \(Q_Du\). It already implies \(\mu(Z)=0\); this is not an
  extra support assumption. Every such measure is the genuine relative
  top product of an actually attained \(u_\mu\in\mathcal E_D^1(L)\),
  unique modulo one real constant. Conversely every finite-energy output
  has (1.7).
\item
  The normalized reference cone is compact; relative envelopes and
  energy upper semicontinuity hold. For every \(u\in\mathcal E_D^1(L)\)
  and every \(f\in C^0(X^{\rm an})\), \[
   \left.\frac{d}{dt}\right|_{0}E_D^\downarrow(u+tf)
                    =\int f\,\mathcal M_D(u,\ldots,u),       \tag{1.8}
  \] with both derivative signs. The normalized variational problem is
  coercive in relative \(J\) and its supremum is attained by an actual
  function. These are proved in Sections 7--9.
\end{enumerate}

The entire positivity domain is the set of all tuples
\((k,X,L,A,D,(1.1))\) just specified. In particular, forgetting the
reference data projects onto \textbf{every} actual big rational line
bundle, including every nonnef and every non-finitely-generated big
bundle. Within each datum (1.3) is the entire genuine dominated cone,
not a selected FS, toric, skeletal or model-supported subcone. The
reference and the raw mass are retained in every definition and in the
equation. Different Kodaira data need not give the same theory. Neither
\(\sigma=V_L\) nor \(V_D=\operatorname{vol}(L)\) is asserted. We do not
treat every possible reference type or every pseudoeffective nonbig
bundle.

\textbf{Theorem 2 (additional zero-volume range).} On the same \(k,X\),
the actual trivial rational bundle \(L=\mathcal O_X\) has
\(\operatorname{PSH}(0)=\mathbb R\), reference zero, raw energy
identically zero, finite-energy space \(\mathbb R\), and genuine top
product zero. The only positive Radon target of the correct mass is the
zero measure. Every actual finite constant solves its equation, and
solutions are unique modulo constants. Compactness, envelopes, both
derivatives, comparison and attainment are elementary. This is exactly
the extra zero-volume equation range claimed. There is no division by a
zero mass.

In dimension zero \(X=\operatorname{Spec}k\) instead has the empty top
product of mass one and energy \(E(a)=a\). It is not a
positive-dimensional zero-volume case. Sections 10--11 verify Theorem 2
and the two diagnostics separately.

\section{2. Mathematical premises and source
discipline}\label{mathematical-premises-and-source-discipline}

\textbf{G0, the task grant.} We use the exact supplied full actual-ample
P16.1 mathematical body, in its stated domain: every reduced irreducible
geometrically unibranch projective variety over every algebraically
closed trivially valued characteristic-zero field, polarized by an
actual ample rational line bundle. Its genuine ample energy and raw
mixed operator, genuine decreasing and weighted limits, envelope
property, normalized compactness, comparison and domination, two-sided
continuous outer derivative, intrinsic measure-energy criterion,
coercivity, target \(L^1\) continuity, attained equation, and uniqueness
modulo constants are available. Its raw mass is the ample intersection
number. Here it is applied only to the actual ample \(A\) on the
original smooth \(X\). No assertion about a nonample bundle is imported
from that grant. The proof is preserved as
\texttt{preserved/block-00.md}; its primary proof dictionary is
preserved separately. Audit labels and supplied hashes are documentary
identifiers, not new premises. The incomplete102 endpoint and
constructions106--107 are not premises.

\textbf{External published inputs.} BJ is \emph{Global pluripotential
theory over a trivially valued field}, arXiv:1801.08229v4. BJ+ is the
addendum arXiv:2206.07183v2, PDF dated 7 September 2023. BE is
\emph{Spaces of norms, determinant of cohomology and Fekete points in
non-Archimedean geometry}, arXiv:1805.01016v3. CLD is arXiv:1204.6277v2.
The exact ledger in Appendix A and \texttt{external-premises.json}
states the scopes and preserves the supplied literal coordinates. Only
the complete pertinent statement/proof units are claimed read, not
entire books or recursively all proofs they cite.

BJ+ Theorem A already proves the envelope property for \textbf{every
numerical class} on smooth projective \(X\) in the retained
characteristic-zero field domain. Its proof uses smooth test
configurations, multiplier-ideal subadditivity, uniform positive-twist
generation, and decreasing divisibility-indexed approximants; divisorial
determination identifies the limit envelope. Corollaries 2.1--2.2
already give the genuine \(V_\theta=P_\theta(0)\) for every
pseudoeffective class and normalized compactness. Corollary 2.3 states
full Siu division. These are established foundations, not new results or
missing nonnef inputs. BJ+ Theorem B still assumes ample \(\omega\) and
supplies no nonample energy theorem by changing that adjective. We
reconstruct the ample-remainder division directly in Section 4,
repairing the particular faulty monotonicity inference in the older
printed proof of BJ 6.21.

The bigness convention for an actual integral bundle \(B\) is the usual
positive section volume: \[
 \limsup_{m\to\infty}\frac{n!\,h^0(X,mB)}{m^n}>0.              \tag{2.1}
\] For rational bundles clear a denominator. No existence of a limit or
finite generation is used below. This is the standard algebraic meaning
of an actual big rational line bundle; (2.1) supplies an infinite
sequence with positive \(m^n\) growth. Section 3 proves the Kodaira
characterization needed here rather than assuming it.

\section{3. Every actual big rational bundle: a proof without
Bertini}\label{every-actual-big-rational-bundle-a-proof-without-bertini}

Clear \(rL\) to an integral big bundle \(B\). Choose an actual very
ample \(M\) and a nonzero section. Since \(X\) is integral, its zero
scheme \(H\) is an effective Cartier divisor (possibly empty). If empty,
\(M\) would be trivial, impossible in positive dimension for a very
ample bundle on a positive-dimensional proper variety. The projective
Cartier scheme \(H\) has dimension \(n-1\); reducedness and
irreducibility of \(H\) are unnecessary.

Choose a sufficiently positive power \(T\) of a very ample bundle on
\(H\) so that \(T\otimes(B|_H)^{-1}\) is globally generated. This is
Stacks 01Q3(7), on the possibly nonreduced projective \(H\). There are
finitely many associated points. At each such point the sections
vanishing in its fiber form a proper \(k\)-linear subspace of the
finite-dimensional global section space: global generation supplies a
nonvanishing section there. A finite union of proper linear subspaces
cannot cover a vector space over infinite \(k\). Choose a section \(s\)
outside their union.

This is a regular section. For completeness, the local zero-divisor
criterion is Stacks 00LD: if a nonzero local element annihilates an
element of a module, its nonzero annihilated submodule has an associated
prime containing that element, which is also associated to the ambient
module. Conversely membership in an associated prime produces an
annihilated nonzero element. Applied on affine trivializations this says
that a section nonvanishing at every associated point is a
nonzerodivisor. This also supplies the omitted argument referenced in
Stacks 0B3Q, Lemma 31.16.1. Thus multiplication by \(s^m\) injects \[
                 B^m|_H\lhook\joinrel\longrightarrow T^m.
\] Serre vanishing on proper \(H\) (Stacks 0B5T(4)) and the degree bound
for its Euler-characteristic polynomial (Stacks 0BEM) give \[
 h^0(H,B^m|_H)\leq h^0(H,T^m)=O(m^{n-1}).                     \tag{3.1}
\] In dimension one this is simply a constant length bound. Bigness
gives \(h^0(X,mB)\geq c m^n\) for infinitely many arbitrarily large
\(m\). The exact Cartier restriction sequence \[
 0\longrightarrow mB-H\longrightarrow mB\longrightarrow mB|_H
                 \longrightarrow0
\] and (3.1) therefore give a nonzero section of \(mB-H\) for some
\(m\). On integral \(X\) its zero divisor \(D'\) is effective Cartier
and it provides an actual bundle isomorphism
\(mB\simeq M\otimes\mathcal O(D')\). Consequently \[
                 A=M/(mr),\qquad D=D'/(mr)
\] give (1.1). Thus the datum does not silently restrict the big input
domain. No Bertini, contraction, section-ring finite generation,
countable field or target measure has entered.

\section{4. The whole genuine reference cone and its actual
inverse}\label{the-whole-genuine-reference-cone-and-its-actual-inverse}

\subsection{4.1. Genuine addition, injectivity and all-point
order}\label{genuine-addition-injectivity-and-all-point-order}

The divisor potential is a generically finite FS potential in the actual
bundle \(\mathcal O(D)\). On a same-bundle FS approximant of an
\(A\)-potential, multiplying its finite sections by the cleared section
\(s_{mD}\) gives its image under (1.4) in the actual \(L\). Those images
decrease with the original approximants and are generically finite.
Genuine psh addition, BJ 4.7, gives the same conclusion without choosing
approximants. Their pointwise sum is defined even on \(Z\), since no
summand takes \(+\infty\). Thus \(F_D\psi\) is a genuine \(L\)-psh
function with value \(-\infty\) at every point of \(Z\).

If \(F_D\psi=F_D\chi\), subtract finite \(\sigma\) on all divisorial
points. BJ 4.23 implies \(\psi=\chi\) everywhere. Similarly, if
\(F_D\psi\leq F_D\chi\), then \(\psi\leq\chi\) divisorially, and hence
everywhere by BJ 4.22. This proves order reflection, including values on
the fixed pole set. It also shows how an inverse, once existence is
proved, must be determined; there is no choice of a value over a divisor
point.

\subsection{4.2. Homogeneous division with a compactness
repair}\label{homogeneous-division-with-a-compactness-repair}

Take homogeneous \(h\in\operatorname{PSH}(A+D)\) with \(h\leq\sigma\).
If \(D=0\) there is nothing to divide. Otherwise BJ 6.17 supplies a
decreasing net of homogeneous FS functions \[
 h_i\in\mathcal H_{\rm hom}(L_i),\quad h_i\downarrow h,
          c_1(L_i)\longrightarrow c_1(A+D),\quad \sup h_i=0.
\] The extended-continuous section potential has a compact level set
\(K=\{\sigma=-1\}\). For each rational \(0<\varepsilon<1\), the
decreasing compact closed sets \[\{x\in K:h_i(x)\geq-1+\varepsilon\}\]
have empty intersection; their complements cover \(K\), and directedness
yields an empty-set tail beyond an index \(i(\varepsilon)\). Use the
directed refinement of valid pairs \((i,\varepsilon)\) with
\(i\geq i(\varepsilon)\), ordered by increasing \(i\) and decreasing
positive rational \(\varepsilon\). A common upper pair is obtained by
choosing a smaller tolerance and an index above the two old indices and
its new \(i(\varepsilon)\). This refinement retains cofinal convergence
in the original index. Put \(t_{i,\varepsilon}=1-\varepsilon\). We
obtain \(t_{i,\varepsilon}>0\) tending to one and
\(h_i\leq t_{i,\varepsilon}\sigma\) on \(K\). At a divisorial \(v\) with
\(\sigma(v)<0\), scale \(v\) to \(K\) and use homogeneity. At
\(\sigma(v)=0\) use \(h_i(v)\leq\sup h_i=0\). BJ 4.22 then proves the
inequality on every point, using the usc obstacle
\(t_{i,\varepsilon}\sigma\).

BJ 2.23 is actual section division on a normal variety. After clearing
denominators the bound forces every FS section to vanish to the
specified divisor order: replace \(v\) by \(av\) and let \(a\to+\infty\)
to remove the additive bound; height-one valuations and normality then
give divisibility by the divisor section. Its quotient formula provides
actual homogeneous FS extensions \[
 r_i=h_i-t_i\sigma\in\mathcal H_{\rm hom}(L_i-t_iD),\qquad
 \theta_i=c_1(L_i)-t_i[D]\longrightarrow c_1(A).               \tag{4.1}
\] The displayed subtraction is initially on \(U\); section division
provides its values on \(Z\).

\textbf{The net \((r_i)\) need not decrease.} For example, at a point
with \(\sigma=-1\), numerical values \(h_1=-0.8\), \(t_1=0.7\) and
\(h_2=-0.9\), \(t_2=0.85\) satisfy both required upper bounds while
\(h_2<h_1\) and \(r_2=-0.05>r_1=-0.1\). This checks the inference, not a
counterexample to the division theorem. We do not use that inference in
BJ 6.21.

Instead use openness of the ample cone in finite-dimensional \(N^1(X)\).
There are rational \(c_i>1\), tending to one, with
\(c_i c_1(A)-\theta_i\) ample. One can choose \(c_i-1\) slightly larger
than a fixed multiple of \(\|\theta_i-c_1(A)\|\), after restricting to a
tail. Adding the zero potential in this ample difference and dividing by
\(c_i\) gives \[
             r_i/c_i\in\operatorname{PSH}(A),\qquad
                       (r_i/c_i)(v_{\rm triv})=0.
\] Normalized ample compactness, G0, yields a subnet converging to an
\textbf{actual} \(r\in\operatorname{PSH}(A)\). At every divisorial
\(v\), all terms are finite and \[
             (r_i/c_i)(v)\longrightarrow h(v)-\sigma(v).
\] Hence \(F_Dr\) and \(h\) agree divisorially, and both are genuine
\(L\)-psh functions. BJ 4.23 gives \(F_Dr=h\) at every analytic point.
Homogeneity of \(r\) follows by the same divisorial determination. This
proves the needed homogeneous division without a monotone divided net.

\subsection{4.3. Arbitrary potentials and the
poles}\label{arbitrary-potentials-and-the-poles}

Let \(u\in\operatorname{PSH}(L)\) satisfy \(u\leq\sigma+C\) and put
\(M=\sup u<\infty\). BJ 6.24 gives homogeneous
\(h_\lambda\in\operatorname{PSH}_{\rm hom}(L)\) for \(\lambda<M\) such
that \[
 h_\lambda+\lambda\leq u,\qquad
              u=\sup_{\lambda<M}(h_\lambda+\lambda)             \tag{4.2}
\] pointwise. For any divisorial \(v\) and any \(a>0\), homogeneity and
domination give \[
 a h_\lambda(v)\leq a\sigma(v)+C-\lambda.
\] Let \(a\to+\infty\), then apply BJ 4.22: \(h_\lambda\leq\sigma\)
everywhere. Section 4.2 provides actual homogeneous
\(r_\lambda\in\operatorname{PSH}(A)\) with \(F_Dr_\lambda=h_\lambda\)
and \(\sup r_\lambda=0\).

The family \(r_\lambda+\lambda\) is uniformly bounded above by \(M\).
Ample EP gives \[
                r=\sup^*_{\lambda<M}(r_\lambda+\lambda)
                                      \in\operatorname{PSH}(A).
\] BJ 5.6 says that usc regularization changes none of the divisorial
supremum values. On those points (4.2) therefore gives \(F_Dr=u\). BJ
4.23 again gives equality everywhere, including all points on \(Z\),
where the image is \(-\infty\). Conversely every
\(\psi\in\operatorname{PSH}(A)\) is bounded above, so
\(F_D\psi\leq\sigma+\sup\psi\). This proves exactly (1.3)--(1.4).

Define \(Q_Du\) as this unique extension of \(u-\sigma\) on \(U\). It is
intrinsic to \(u\) and the retained actual datum. Its divisorial
coordinates are \(u(v)-\sigma(v)\); adding/subtracting these fixed
finite coordinates proves continuity of both maps. The maximum principle
gives \[
 \sup u=u(v_{\rm triv})=Q_Du(v_{\rm triv})=\sup Q_Du.          \tag{4.3}
\] Thus \(\sigma=F_D0\) is the maximal normalized member of the
\textbf{relative} cone. It is generally different from the already
established unrestricted \(V_L\).

If \(u_i\downarrow u\) genuinely in the relative cone, then
\(Q_Du_i\geq Q_Du_j\) for \(i\leq j\) by order reflection. The infimum
has the prescribed finite divisorial and generic values; decreasing psh
closure and determination identify it with \(Q_Du\) at every point. This
fact will justify all unbounded decreasing transports, not merely weak
divisorial convergence.

\section{5. Genuine local metric products, constructed before a
target}\label{genuine-local-metric-products-constructed-before-a-target}

\subsection{5.1. The signed frame
calculation}\label{the-signed-frame-calculation}

First clear (1.1) so frames are honest. On a trivializing algebraic
open, take \(e_A\) and \(e_D\) with \(s_D=f e_D\), and
\(e_L=e_A\otimes e_D\). Write potentials using \[
 \|e\|_\psi=e^{-\psi}\|e\|_{\rm triv},\qquad
                 h_A=-\log\|e_A\|_\psi.
\] Since \(\sigma=\log|f|+\log\|e_D\|_{\rm triv}\) and
\(u=\sigma+\psi\), the actual local weight on \(e_L\) is \[
 h_L=u-\log\|e_A\|_{\rm triv}-\log\|e_D\|_{\rm triv}
                         =h_A+\log|f|.                        \tag{5.1}
\] The sign is \textbf{plus}. In the frame \(e_A\otimes s_D=f e_L\) on
\(U\) it becomes \(h_L-\log|f|=h_A\). In particular \(\|s_D\|\) for the
divisor metric equals one on \(U\), and multiplication by the actual
\(s_D\) identifies the metrized bundles \(A\) and \(L\) there. Rational
cases divide this identity by the common denominator.

On \(U\), \(f\) is an analytic unit. CLD 18.1.2 gives \(dd^c\log|f|=0\)
and the usual curvature convention \(dd^c(-\log\|e\|)\). Its local
Poincare--Lelong proof also gives this equality. CLD 8.8.5 proves
harmonic-addition invariance of continuous local Bedford--Taylor
products; it does not require introducing an inverse target. Therefore
(5.1) identifies smooth local products on the \textbf{original} \(U\),
with the original analytic fundamental cycle.

\subsection{5.2. Continuous and bounded
stages}\label{continuous-and-bounded-stages}

For continuous genuine \(A\)-psh metrics, BE 7.14 provides decreasing
smooth semipositive regularization on the actual semiample \(A\),
including trivial fields. On compact \(X^{\rm an}\) a continuous limit
gives uniform convergence. BE 8.3--8.5 construct their genuine mixed
measures and their unique uniform-to-weak extension, preserving
symmetry, positivity and raw intersection mass. At FS/model stages their
cycle coefficients are actual intersection coefficients; BE 8.18 and BJ
7.23 identify the continuous local product with the genuine ample
energy-pairing operator. Thus the continuous product identity on \(U\)
holds before singular extensions: \[
 \prod_i c_1(L,u_i)|_U
        =\left.\bigwedge_i(c_1(A)+dd^c\psi_i)\right|_U,
                     \qquad u_i=F_D\psi_i.                   \tag{5.2}
\] We do not use BE 7.8, which assumes a nontrivial valuation and an
ample bundle. Nor do we claim global continuous semipositive
regularization on nonnef \(L\): only local transport of actual \(A\)
metrics on \(U\) is used.

Next let the quotients \(\psi_i\) be bounded genuine psh functions.
Choose their genuine decreasing continuous same-class approximants
supplied by G0. These give decreasing genuine \(L\) metrics, locally
bounded on \(U\): the finitely many frame/reference functions are
bounded on each compact subset there. Define the local extended product
by their local weak limits. For any \(g\in C_c(U)\) its zero extension
is continuous on compact \(X^{\rm an}\), so G0 decreasing continuity and
(5.2) show that its limit integral is \[
                    \int_U g\,\bigwedge_i(A+dd^c\psi_i).      \tag{5.3}
\] This proves existence and independence of all chosen genuine
decreasing original-reference approximants. It also proves frame
independence via (5.1) and harmonic invariance. The product is local on
ordinary open sets by genuine ample mixed locality. It is a construction
on this specified family of local singular metrics; we do not assume a
universal local product for arbitrary nonample metrics. Every bounded
stage comes from genuine local smooth/continuous metric products, and
the limit is determined by those products, not by a selected target.
This is the required local finite-energy extension in this cone.

The original intersection normalization is fixed: at an ample model
stage the raw measure is \[
 \sum_E b_E\bigl((A_{\mathcal X}+D_1)|_E\cdots
                         (A_{\mathcal X}+D_n)|_E\bigr)\delta_{v_E},
                \qquad b_E=\operatorname{ord}_E(\mathcal X_0). \tag{5.4}
\] The relative metric identity uses that same cycle on \(U\). No
component multiplicity, degree, carrier or probability factor is
changed.

\subsection{5.3. Zero pluripolar mass}\label{zero-pluripolar-mass}

BJ 4.36 proves that every proper algebraic closed subset, including
\(Z\), is pluripolar. BJ 5.3 proves that pluripolar Borel sets are
negligible. BJ 13.2 proves that finite-dual-energy measures do not
charge negligible Borel sets, under ample EP; its proof uses genuine
increasing psh approximants, divisorial nonnegligibility, target \(L^1\)
continuity and monotone integration, and does not use an inverse
equation. By G0 every finite-energy ample output has finite dual energy.
Hence every finite-energy top product vanishes on pluripolar sets.

For a mixed product \(M(\psi_1,\ldots,\psi_n)\), let
\(\bar\psi=n^{-1}\sum_i\psi_i\in\mathcal E^1(A)\). Positivity and
multilinearity give the measure inequality \[
  M(\psi_1,\ldots,\psi_n)
       \leq \frac{n^n}{n!}(A+dd^c\bar\psi)^n.                 \tag{5.5}
\] Indeed the expansion of the right top product contains the mixed term
\(n!\,n^{-n}M\) and all other terms are positive. Thus mixed
finite-energy products also vanish on every pluripolar Borel set. In
particular (5.3) has zero mass on the removed fixed divisor in its
global ample extension.

\subsection{5.4. Nonpluripolar reference
truncation}\label{nonpluripolar-reference-truncation}

For arbitrary \(u\in\mathcal P_D(L)\) put \(\psi=Q_Du\) and use the
actual truncations \[
 \psi_j=\max(\psi,-j),\qquad
 u_j=F_D\psi_j=\max(u,\sigma-j),\qquad
                 B_j=U\cap\{\psi>-j\}.                      \tag{5.6}
\] The maximum identity also holds at \(Z\) with the extended-value
convention. These truncations are genuine relative \(L\) metrics with
bounded quotients. We never use \(\max(u,-j)\), since the constant floor
is not \(L\)-psh for nonnef \(L\).

Let \(\nu_j\) be the zero extension of the bounded local top product of
\(u_j\), restricted to \(B_j\). For \(k\geq j\),
\(\psi_j=\max(\psi_k,-j)\) and \(\{\psi_k>-j\}=\{\psi>-j\}\). Genuine
strict-winning locality at the bounded stages gives \[
 \mathbf1_{B_j}(A+dd^c\psi_j)^n
                 =\mathbf1_{B_j}(A+dd^c\psi_k)^n.
\] Consequently \(\nu_j\leq\nu_k\) as measures and
\(\nu_j(X^{\rm an})\leq V_D\). Their increasing limit is a positive
Radon measure: for a Borel set, increasing convergence and inner
regularity at a sufficiently large finite stage give inner regularity of
the limit; bounded mass then gives the usual finite regular measure on
compact \(X^{\rm an}\). This defines the genuine \textbf{relative
nonpluripolar top product} by local metric products and reference
truncations, before any target is considered. Replacing \(j\) by any
increasing floor levels tending to infinity gives the same limit, by the
same locality and interlacing argument.

If \(\psi\in\mathcal E^1(A)\), locality applies directly between
\(\psi\) and \(\psi_j\) and gives \[
       \nu_j=\mathbf1_{B_j}(A+dd^c\psi)^n.                    \tag{5.7}
\] The output integrates \(\psi\), hence charges none of
\(\{\psi=-\infty\}\); it also vanishes on \(Z\) by Section 5.3. The
\(B_j\) therefore exhaust a set of full output measure. Equations (5.7)
and monotone integration prove \[
             \lim_j\nu_j=(A+dd^c\psi)^n,
                    \qquad \text{raw mass }V_D.               \tag{5.8}
\] This is not merely a weak limit whose mass may disappear at the fixed
pole. It is an increasing nonpluripolar restriction identity with no
lost mass.

For mixed finite-energy products the bounded local product is already
symmetric and multilinear. One may equivalently use the positive
polarization identity \[
 M(\psi_1,\ldots,\psi_n)=\frac1{n!}
  \sum_{S\subset\{1,\ldots,n\}}(-1)^{n-|S|}|S|^n
     \operatorname{MA}_{\rm raw,A}\left(|S|^{-1}\sum_{i\in S}\psi_i\right), \tag{5.9}
\] where the empty term is zero. Formula (5.9) follows by
inclusion--exclusion applied to the multilinear expansion. Each average
is genuine finite energy. Its local nonpluripolar top product has just
been proved equal to the genuine ample top product. The polarization
therefore equals the positive mixed product in (1.6), is local, and
vanishes on all the factors' pole sets by (5.5) and their integrability.
Alternatively use simultaneous bounded truncations and the mixed
locality argument. Both constructions have the same local continuous
starting product and the same finite-energy limit. Nothing in this
construction refers to an equation target.

\section{6. Actual relative energy and all decreasing
limits}\label{actual-relative-energy-and-all-decreasing-limits}

For bounded quotients, define the actual relative energy by \[
 E_D(u)=\frac1{n+1}\sum_{j=0}^n
       \int_{X^{\rm an}}Q_Du\,
          \mathcal M_D(\underbrace{u,\ldots,u}_{j},
                         \underbrace{\sigma,\ldots,\sigma}_{n-j}). \tag{6.1}
\] On \(U\) the integrand is the actual relative difference
\(u-\sigma\). Its uniquely determined extension \(Q_Du\) makes it
meaningful on all points without subtracting two infinite values.
Section 5 identifies every measure in (6.1) with the genuine ample mixed
metric product. The ample energy identity consequently proves \[
                     E_D(u)=E_{{\rm raw},A}(Q_Du),\qquad E_D(\sigma)=0.
\] For arbitrary \(u\) use (5.6) and define \(E_D(u)=\lim_j E_D(u_j)\)
in \(\mathbb R\cup\{-\infty\}\). Monotonicity gives the limit, and
genuine ample decreasing energy continuity proves the same identity for
every actual quotient. Thus (1.5) is an identity between intrinsically
defined actual spaces, not an abstract completion or an arbitrary lift
on a modification. For a finite-energy quotient the mixed integrals in
(6.1) remain finite and give the same value.

Raw energy translates by \(V_D c\): \[
             E_D(u+c)=E_D(u)+V_Dc.                            \tag{6.2}
\] It is increasing, concave, weakly usc and decreasing-net continuous
in its genuine domain. All follow from the actual quotient identity and
G0, with order and all-point decreasing convergence verified in Section
4.3. Relative bounded potentials lie in \(\mathcal E_D^1\), and the
domain is upward closed within the genuine reference cone.

For genuine \(u_{ij}\downarrow u_i\) with finite-energy endpoints,
Section 4.3 gives genuine \(Q_Du_{ij}\downarrow Q_Du_i\) at every point.
G0 mixed decreasing continuity then proves \[
                \mathcal M_D(u_{1j},\ldots,u_{nj})
                        \longrightarrow\mathcal M_D(u_1,\ldots,u_n)
\] weakly against every continuous test on the original compact carrier.
These may be unbounded nets over an arbitrarily large field. Their
masses remain exactly \(V_D\). For completeness the proof tests against
real PL functions using the decreasing signed energy pairing, then uses
uniform PL density and the fixed positive masses for general continuous
tests.

For a genuine quotient test \(\psi_0\in\operatorname{PSH}(A)\)
integrable against the limiting mixed measure, and decreasing genuine
approximants \(\psi_{0j}\downarrow\psi_0\), the weighted limit is \[
 \psi_{0j}\,\mathcal M_D(u_{1j},\ldots,u_{nj})
                     \longrightarrow \psi_0\,\mathcal M_D(u_1,\ldots,u_n). \tag{6.3}
\] The correctly signed total integral is \[
 (c_1(A),\psi_0)\prod_{i=1}^n(c_1(A),Q_Du_i)
          -(c_1(A),0)\prod_{i=1}^n(c_1(A),Q_Du_i).             \tag{6.4}
\] Decreasing pairing continuity gives convergence of these totals.
Translate the usc tests below zero; their negatives are increasing lsc
functions. For positive weakly convergent measures the lsc integral
inequality implies that any weighted subsequential limit dominates the
claimed limit. Equality of the limiting totals proves equality. This is
the ample weighted argument in G0/BJ 7.21(ii), with (6.4) correcting the
reversed order in its printed proof. A difference of quotient tests is
handled by two converging terms, not by claiming that their difference
decreases. The word weighted here means \textbf{reference-relative
quotient weight}. Absolute integrability of \(\sigma\) or \(u\) against
every finite-energy measure is not inferred.

\section{7. Relative envelopes, compactness and all original
directions}\label{relative-envelopes-compactness-and-all-original-directions}

Equation (4.3) and the divisorial homeomorphism identify the normalized
relative cone with normalized \(\operatorname{PSH}(A)\). G0 gives
compactness, even when the original field and weak space are
nonmetrizable. In particular a compactness limit is an actual genuine
relative function \(F_D\psi\) on all original points. Its specified
reference is retained.

For \(b\in C^0(X^{\rm an})\) define the relative envelope as the
supremum in (1.3) below \(\sigma+b\). For an actual
\(\chi\in\operatorname{PSH}(A)\), \[
 F_D\chi\leq\sigma+b
     \quad\Longleftrightarrow\quad \chi\leq b.                \tag{7.1}
\] The forward direction subtracts only on \(U\), hence holds
divisorially, then BJ 4.22 extends it to every original point because
\(b\) is usc. The reverse direction is immediate from addition,
including at \(Z\). Therefore \[
                    P_D(\sigma+b)=F_D(P_A(b)).                \tag{7.2}
\] Ample EP supplies a continuous actual quotient \(P_A(b)\). More
generally \(\sigma+w\) with upper bounded usc \(w\) has the same
envelope identity whenever its ample envelope is genuine. Again (7.1) is
proved by divisorial domination, not by subtraction at a pole. The same
argument gives \[
                 P_D(u+tf)=F_D(P_A(Q_Du+tf))                  \tag{7.3}
\] for every real \(t\) and \(f\in C^0(X^{\rm an})\).

Define the relative outer energy by genuine relative continuous
majorants: \[
 E_D^\downarrow(\sigma+w)
   :=\inf_{b\in C^0,\ \sigma+b\geq\sigma+w}E_D(P_D(\sigma+b)). \tag{7.4}
\] For \(w\) upper bounded usc, the majorant condition is equivalent to
\(b\geq w\) whenever \(w\) is the quotient obstacle used here
(\(w=\psi+tf\), or continuous \(w\)); the necessary extension follows as
in (7.1), applied to \(\psi\leq b-tf\). Thus (7.4) equals
\(E_{{\rm raw},A}^{\downarrow}(w)\) for all the asserted obstacles. No
undefined quotient of a non-psh object at \(Z\) is used: \(w\) is the
specified actual quotient obstacle on the whole carrier. G0 outer energy
agrees with energy at psh functions and with the envelope energy under
EP.

For a real PL \(f\), the actual ample derivative estimate gives \[
 E_D^\downarrow(u+tf)=E_D(u)+t\int f\,\mathcal M_D(u^n)+O_f(t^2), \tag{7.5}
\] uniformly in \(u\in\mathcal E_D^1\). Negative \(t\) follows by
\(-f\). The unbounded version uses genuine decreasing continuous
quotients and G0 outer-energy and mixed continuity, so its uniform error
survives the passage to the endpoint.

For arbitrary original \(f\in C^0\), choose real PL \(f_m\) with
\(\|f-f_m\|_\infty\leq\varepsilon_m\to0\). For all real \(t\), \[
 u+tf_m-|t|\varepsilon_m\leq u+tf\leq u+tf_m+|t|\varepsilon_m.
\] Outer monotonicity and raw translation bound the difference of
energies by \(V_D|t|\varepsilon_m\). First divide by positive or
negative \(t\), with the inequality reversed in the latter case, and use
(7.5); then let \(m\to\infty\). Both one-sided limits equal the integral
in (1.8). The error is \(|t|\varepsilon_m\), not \(t^2\varepsilon_m\),
and the upper derivative error has a plus sign. The direction is any
continuous function on the original \(X^{\rm an}\); no pullback or
retraction restricts it.

\section{8. Comparison, locality, domination and
uniqueness}\label{comparison-locality-domination-and-uniqueness}

Let \(u=F_D\psi\), \(v=F_D\chi\) be genuine relative finite-energy
potentials. Their strict winning sets on \(U\) are precisely those of
\(\psi,\chi\). At \(Z\) both image functions are \(-\infty\), so neither
strict inequality holds. All their relevant products vanish on \(Z\).
Since \(\max(u,v)=F_D\max(\psi,\chi)\), genuine ample strict-winning
locality proves \[
 \mathbf1_{\{u>v\}}\mathcal M_D(\max(u,v)^n)
                    =\mathbf1_{\{u>v\}}\mathcal M_D(u^n).
\] These Borel sets need not be ordinary open sets. At the bounded
stages this comes from the model coefficient formula (5.4); unbounded
passage uses quotient-weighted continuity term by term. Mixed locality
on ordinary open sets follows by polarization or directly from ample
mixed locality. Reference zero extension is essential in each equality.

Displaced maxima and the equal raw mass \(V_D\) give the actual
comparison principle \[
    \int_{\{u<v\}}\mathcal M_D(v^n)
                  \leq\int_{\{u<v\}}\mathcal M_D(u^n).         \tag{8.1}
\] Equivalently apply the ample comparison inequality on
\(\{\psi<\chi\}\) and delete its zero measure part in \(Z\). This proof
uses no inverse equation. If \(a\in\mathcal P_D(L)\),
\(b\in\mathcal E_D^1\) and \(a\leq b\) almost everywhere for
\(\mathcal M_D(b^n)\), then quotient inequality holds almost everywhere
off \(Z\). G0 ample domination yields \(Q_Da\leq Q_Db\) everywhere;
addition proves actual domination on all original points.

Use the probability calibration only when \(V_D>0\) to set \[
 I_D(u,v)=V_D^{-1}\int(Q_Du-Q_Dv)
                         (\mathcal M_D(v^n)-\mathcal M_D(u^n)). \tag{8.2}
\] Integrability of the quotients makes this well defined, and it equals
the genuine ample \(I_A\). Thus \(I_D\geq0\), and \(I_D=0\) forces
\(Q_Du-Q_Dv\) to be one real constant by ample uniqueness on integral
\(X\). Consequently \(u=v+c\) everywhere, with \(-\infty+c=-\infty\) at
the fixed poles. Equal genuine relative measures imply this equality
modulo constants. Constants plainly leave products unchanged. This is
actual metric uniqueness, not uniqueness of an arbitrary lift.

\section{9. Intrinsic targets, coercivity and actual
attainment}\label{intrinsic-targets-coercivity-and-actual-attainment}

Put \(\bar\mu=\mu/V_D\). From (1.5), the onto actual cone identity, and
raw calibration, \[
 D_D(\mu)=E_{{\rm raw},A}^{\vee}(\mu)
                      =V_D E_A^{\vee}(\bar\mu).               \tag{9.1}
\] No support hypothesis is hidden in this equality. G0 says that its
finiteness is equivalent to integrability of \textbf{every actual}
\(\mathcal E^1(A)\) potential, hence every quotient in (1.7). The proof
of that criterion forms actual decreasing convex sums in a contradiction
argument on a bounded \(J\) set, then uses the quadratic small-scaling
\(J\) estimate to bound the dual supremum. This proof concerns the
original function space and does not posit an inverse target. BJ 13.2
applied to \(\bar\mu\) implies \(\mu(Z)=0\). Requiring that equality
separately would be redundant. Neither absolute integration of \(u\) nor
integration of \(\sigma\) substitutes for the quotient integral.

For an explicit attainment proof, write \(d=E_A^{\vee}(\bar\mu)<\infty\)
and normalize \(\sup u=0\). Define \[
 J_D(u)=\sup Q_Du-V_D^{-1}E_D(u),\qquad
 J_{D,\mu}(u)=d-V_D^{-1}E_D(u)+\int Q_Du\,d\bar\mu.
\] They are the actual ample functionals on the genuine quotient, and G0
strict concavity gives \[
                   J_D(u)\leq C_n(J_{D,\mu}(u)+d).            \tag{9.2}
\] A normalized actual maximizing sequence for (1.7) has eventually
bounded \(J_D\). Section 7 compactness gives a subnet converging to an
actual \(u=F_D\psi\) in the normalized genuine relative cone. Weak lower
semicontinuity of \(J_D\) makes \(E_D(u)>-\infty\). G0 target \(L^1\)
continuity on bounded-\(J\) sets, applied to the actual quotients, gives
convergence of their \(\mu\) integrals. Weak energy upper semicontinuity
now proves \[
 E_D(u)-\int Q_Du\,d\mu
       \geq\limsup_i\left(E_D(u_i)-\int Q_Du_i\,d\mu\right)=D_D(\mu).
\] The reverse inequality is the definition of the supremum. Thus it is
actually attained. There is no metrizability assumption on the compact
space: a maximizing sequence may have a convergent subnet. The entire
reference cone identity ensures that this compactness step does not lose
the prescribed singularity type.

To verify the equation from the attained maximizer, fix
\(f\in C^0(X^{\rm an})\). The ample envelope of \(\psi+tf\) lies between
\(\psi\pm|t|\|f\|_\infty\), is genuine finite energy, and (7.3) retains
the original reference. Duality for its quotient and monotonicity of
integration give, for every real \(t\), \[
 E_D^\downarrow(u+tf)-\int(\psi+tf)\,d\mu
                   \leq E_D(u)-\int\psi\,d\mu.              \tag{9.3}
\] The first term equals the energy of that actual envelope; since the
envelope lies below the obstacle, its own objective bounds the displayed
objective and is at most the attained supremum. Subtract equality at
zero, divide (9.3) by \(t>0\), and use (1.8):
\[\int f\,\mathcal M_D(u^n)\leq\int f\,d\mu.\] Divide by \(t<0\) with
the reversed sign to obtain the opposite inequality. Hence equality
holds for every original continuous direction and the Radon measures
coincide. Alternatively dense PL tests with their common raw mass
suffice. This derivation uses the operator constructed in Section 5, not
an inverse-assigned measure. Section 8 gives uniqueness modulo one real
constant.

Conversely, if \(\mathcal M_D(u^n)=\mu\) with \(u\in\mathcal E_D^1\),
(1.6) identifies \(\mu\) with the actual ample output of \(Q_Du\); G0
energy concavity/duality supplies finite dual energy and its attained
value. Therefore (1.7) is both necessary and sufficient. All admissible
measures have arbitrary support subject to this intrinsic energy
condition. No fixed skeleton, divisorial-support set, torus invariance
or finite-generation condition enters.

\section{10. The exact additional zero-volume
equation}\label{the-exact-additional-zero-volume-equation}

For connected original \(X\), BJ 4.24 gives
\(\operatorname{PSH}(0)=\mathbb R\). Its proof is short: for normalized
\(\varphi\in\operatorname{PSH}(0)\), the decreasing genuine sequence
\(m\varphi\) has generic value zero and a genuine zero-class limit.
Finiteness at every divisorial point forces \(\varphi(v)=0\) there;
divisorial determination gives \(\varphi=0\) everywhere. Translation
gives the constants.

For the actual trivial bundle the genuine continuous metric product
exists by BE 7.14 and 8.3--8.5. In positive dimension its curvature
factors have zero top measure and raw intersection mass zero.
Equivalently constant local weights have zero curvature. The actual raw
energy pairing is zero (every term involves the zero class or a zero
current); define \(E_0(c)=0\) for finite constants and
\(\mathcal E_0^1=\mathbb R\).

Every positive Radon measure of total mass zero is the zero measure. Its
intrinsic dual supremum is \(\sup_{c\in\mathbb R}(0-\int c\,0)=0\) and
is attained by every constant. All solutions differ by constants. The
normalized space is the singleton zero; the envelope of a continuous
obstacle is its infimum constant. The outer energy is zero, so every
original continuous direction has both derivatives zero, equal to its
integral against the zero measure. Comparison, locality and attainment
follow immediately. No coercivity assertion based on positive-volume
normalization is made. This proves exactly Theorem 2; it does not prove
an equation for arbitrary numerically trivial or arbitrary zero-volume
pseudoeffective actual bundles.

\section{11. Two distinct diagnostics and reference
dependence}\label{two-distinct-diagnostics-and-reference-dependence}

\subsection{11.1. Big volume with nonpositive
self-intersection}\label{big-volume-with-nonpositive-self-intersection}

Let \(X=\operatorname{Bl}_p\mathbb P^2\), \(H=\pi^*\mathcal O(1)\) and
\(E\) the exceptional divisor. For rational \(a>0\) put \(L_a=H+aE\).
Here \(H^2=1\), \(H\cdot E=0\) and \(E^2=-1\), the last from its normal
bundle \(\mathcal O_E(-1)\). Thus \[
                     L_a^2=1-a^2,\qquad L_a\cdot E=-a.       \tag{11.1}
\] For divisible \(m\) and each positive integer \(j\leq ma\),
restriction of \(mH+jE\) to \(E\simeq\mathbb P^1\) is
\(\mathcal O_E(-j)\) and has no sections. The successive restriction
sequences therefore give \[
 h^0(X,mL_a)=h^0(X,mH)=\binom{m+2}{2}.                        \tag{11.2}
\] The second equality follows directly from
\(\pi_*\mathcal O_X=\mathcal O_{\mathbb P^2}\): on the affine plane
chart at \(p\), the two blowup charts have coordinate rings \(k[x,y/x]\)
and \(k[y,x/y]\). In their common Laurent fraction ring their
intersection is \(k[x,y]\), as is seen by comparing the unique monomial
exponents: the first chart requires the \(y\) exponent nonnegative, the
second the \(x\) exponent nonnegative, so intersection terms have both
nonnegative. Away from \(p\) the morphism is an isomorphism. Local
trivializations of \(\mathcal O(m)\) then give the asserted equality of
sections. Consequently \(\operatorname{vol}(L_a)=1\) for all \(a>0\). At
\(a=1\) self-intersection is zero; at \(a=2\) it is \(-3\). Both have
\textbf{positive volume} and are nonnef. They are not examples of
zero-volume normalization.

\subsection{11.2. Genuine zero-volume
noninjectivity}\label{genuine-zero-volume-noninjectivity}

Take \(X=\mathbb P^1\times\mathbb P^1\) and \(L=p_1^*\mathcal O(1)\). It
is actual semiample and nef, with \(h^0(X,mL)=m+1\),
\(\operatorname{vol}(L)=0\) and \(L^2=0\). Let \(s_0,s_1\) be the two
pulled homogeneous coordinates; their maximum canonical norm is one
everywhere. For real \(b>0\) the genuine continuous FS potential \[
                  u_b=\max\{\log|s_0|,\log|s_1|+b\}-b       \tag{11.3}
\] has generic trivial value zero. At
\(v=t\operatorname{ord}_{\{s_1=0\}}\), \[
                         u_b(v)=\max(0,b-t)-b.                \tag{11.4}
\] It is negative for \(t>0\) and equals \(-b\) for \(t\geq b\), whereas
the normalized zero potential is zero. Thus these potentials are not
equal modulo constants. BE 7.14 applies to this actual semiample \(L\)
even over the trivial field; BE 8.5 gives its genuine positive
continuous top measure of raw mass \(L^2=0\). Every one of these
measures is therefore zero. This is genuine zero-volume top-operator
noninjectivity, not an arbitrary inverse measure and not the
positive-volume example of Section 11.1. It rules out carrying ample
top-operator uniqueness unchanged to all zero-volume classes. It does
not contradict the trivial-bundle range, whose entire genuine cone is
just the constants.

\subsection{11.3. What is intrinsic and what is
retained}\label{what-is-intrinsic-and-what-is-retained}

An automorphism of the fixed actual rational-bundle identification,
after clearing, is multiplication by a global unit on integral proper
\(X\), hence by \(k^\times\); its trivial norm is one. Likewise another
generator of the same effective divisor changes no canonical potential.
\(Q_D\), the relative integral energy, product and target condition are
therefore intrinsic to the specified actual reference datum and original
function. Changing \(D\) changes that datum. We neither quotient away
such changes nor assert equality with the unrestricted
minimal-singularity cone.

\section{12. Exact scope assessment against original
P16.5}\label{exact-scope-assessment-against-original-p16.5}

The unchanged completion criteria expressly allow a prescribed genuinely
retained reference and a proved relative mass, and ask separately which
extra zero-volume range is established. Under that precise alternative,
Theorem 1 treats the full big rational-line-bundle domain on the
retained smooth projective original \(X\), for every retained field. Its
quantifiers cover every actual big \(L\), every actual Kodaira datum,
every genuine metric dominated by that reference with finite actual
relative energy, and every positive Radon target of the raw relative
mass satisfying its intrinsic energy condition. Theorem 2 supplies the
separately stated additional zero-volume trivial-bundle range.

This is more than a formal axiom transfer or an ample equation with the
word ample removed. Section 3 realizes the geometric reference on every
original big bundle without finite generation. Section 4 proves the
whole genuine cone identity, repairs division, and uniquely recovers
actual values at every point. Section 5 identifies the genuine local
\(L\) metric products with the correctly signed unit-log identity before
any target, then constructs their nonpluripolar extension and proves
mass and truncation identities. Section 6 defines energy by relative
product integrals on actual functions. Sections 7--9 prove the
original-direction envelope inequalities, both signs, compactness,
coercivity, attained maximizer, comparison and uniqueness. All these
arguments are on the original carrier. The ample grant is never applied
to nonample \(L\).

The resulting full big \textbf{prescribed-reference} theorem is claimed
complete relative to G0 and the explicitly scoped published premises. No
mathematical gate in that theorem is left as a formal hypothesis. The
manuscript and canonical result remain unverified drafts. The original
family outside this retained theorem is not declared solved: a
class-only \(V_L\) theory of mass \(\operatorname{vol}(L)\), arbitrary
singularity references, positive parts which are merely movable, and
arbitrary nonbig pseudoeffective equations are separate unproved ranges
here. Full zero-volume uniqueness is in fact incompatible with Section
11.2 without changing the domain or equivalence. Those exclusions are
not hidden restrictions replacing the claimed entire big domain; the
reference is part of the datum permitted by the task.

\section{Appendix A. Exact external-premise
ledger}\label{appendix-a.-exact-external-premise-ledger}

The machine-readable ledger has consecutive IDs E0--E10 and one distinct
grant G0. A record supplies a standalone scope, exact public
version/URL, complete pertinent preserved statement/proof artifacts,
their computed hashes and local one-based inclusive LF coordinates, and
the \textbf{literal supplied} source LF or renderer coordinates where
available. Source-wide hashes supplied in mathematical data remain
supplied bindings unless the actual full bytes were preserved. No
binary-source hash is fabricated. Hashes of locally preserved literal
bodies establish only byte identity.

\begin{itemize}
\tightlist
\item
  \textbf{E0:} BJ Definitions 2.18, 4.1 and Theorems 4.5, 4.7, 4.15,
  with fixed ample-bundle conversion under G0: genuine numerical psh
  cones, section addition, constants in nef classes, positive scaling,
  maxima and generically finite decreasing varying-class closure on
  projective varieties over algebraically closed trivial fields.
  Retained exact units include the flattening proof of 4.5 and the
  entire pertinent 4.15 proof.
\item
  \textbf{E1:} BJ 4.12 and 4.22--4.25: maximum at the generic trivial
  point, divisorial finiteness/determination and extension of
  inequalities to usc obstacles; connected zero-class cone consists of
  constants. Complete proof units, including the short 4.24 proof, are
  preserved.
\item
  \textbf{E2:} BJ 6.15--6.17 and 6.23--6.25: homogeneous FS
  approximation in varying rational bundles, and actual pointwise
  homogeneous decomposition for every genuine psh function. The full
  Legendre-duality argument and finite rational monomial construction
  are preserved. We do not use the faulty divided-net step of 6.21.
\item
  \textbf{E3:} BJ 2.23: on a normal projective variety, a generically
  finite FS function dominated by an effective rational Cartier divisor
  potential is the divisor potential plus an actual FS function in the
  subtracted bundle. Its full height-one section-division proof is
  preserved.
\item
  \textbf{E4:} BJ 5.6 and 5.11: bounded-above psh supremum
  regularization changes no divisorial value; EP is equivalent to
  normalized compactness in the divisorial topology. Pertinent full
  proofs are preserved. Ample compactness itself is also G0.
\item
  \textbf{E5:} BJ+ Theorem A and Corollaries 2.1--2.3, with their
  retained standing hypotheses and complete pertinent proof: smooth
  irreducible projective \(X\), algebraically closed trivial \(k\),
  characteristic zero or dimension at most two; full numerical EP,
  genuine \(V_\theta\) for every psef class, normalized compactness and
  Siu division. Here only characteristic zero is retained; Section 4
  independently repairs the ample-remainder application. Theorem B
  remains ample and is not a nonample premise.
\item
  \textbf{E6:} BJ 4.36, 5.3 and 13.2: proper algebraic closed subsets
  are pluripolar, pluripolar sets negligible, and finite-dual-energy
  measures give zero mass to negligible Borel sets under the standing
  ample EP. The complete pertinent proofs are preserved. Output finite
  dual energy comes from G0.
\item
  \textbf{E7:} BE 7.14, 8.3--8.5, 8.18 and BJ 7.23: on geometrically
  reduced projective schemes over complete non-Archimedean fields,
  semiample rational psh metrics have decreasing smooth semipositive
  regularization, including trivial valuations; smooth mixed measures on
  projective schemes use the original fundamental cycle and raw
  intersection mass; continuous psh-regularizable mixed products have a
  unique uniform-to-weak extension, with model/Deligne compatibility
  identifying the BJ operator. Our smooth \(X\) satisfies all the
  standing geometric hypotheses. The full regularization and
  uniform-extension proofs are retained. BE 7.8 is explicitly excluded.
\item
  \textbf{E8:} CLD 18.1.2, 18.1.6 and 8.8.5: local unit logarithms are
  smooth harmonic, curvature of a metrized line bundle uses
  \(dd^c(-\log\|e\|)\), and continuous local products are invariant
  under harmonic additions. The pertinent local frame and
  Poincare--Lelong proof bodies and product identities are preserved. We
  do not infer arbitrary unbounded local theory from a
  uniform-convergence theorem.
\item
  \textbf{E9:} Stacks 01Q3(7), 0B5T(4) and 0BEM: globally generated
  positive twists on quasi-compact quasi-separated schemes with an ample
  invertible sheaf; Serre vanishing on proper Noetherian schemes;
  numerical Euler polynomial degree bounded by coherent support
  dimension on proper schemes over a field. Full supplied mathematical
  statement/proof bodies are preserved, with their current tag bindings.
  These apply to nonreduced \(H\).
\item
  \textbf{E10:} Stacks 00LD (Lemma 10.63.9), together with 0B3Q (Lemma
  31.16.1): zero divisors on a module over a Noetherian ring are
  precisely the union of associated primes; a line-bundle section is
  regular iff it does not vanish at associated points. The 00LD proof
  was retrieved from the public primary page, including the literal
  renderer coordinates. The 0B3Q proof is printed as omitted with a
  hint; Section 3 supplies its reduction instead of claiming an absent
  proof was read. Finiteness of associated points for a coherent sheaf
  on a Noetherian scheme is the standard associated-prime finiteness
  fact; the public supporting primary body is included in this record.
\end{itemize}

All ample estimates and unbounded pairings used in Sections 6--9 are in
G0 in their actual ample domain. The primary proof dictionary includes
BJ 7.1--7.21, 8.3--8.8, 9.14--9.20, locality/comparison, uniqueness and
12.8, so their precise published proof scopes can be checked without
treating G0 as a nonample premise.

\section{Appendix B. Reproducibility and evidence
limits}\label{appendix-b.-reproducibility-and-evidence-limits}

Run \texttt{python3\ rerun.py} in the worker directory. The script
preserves \texttt{input.json} byte identically, extracts the exact task
and assignment, retains all eight inline mathematical blocks, extracts
literal primary bodies without following any supplied absolute path,
verifies local hashes and original coordinate metadata, checks the
premise and backward-support ledgers, reruns the exact rational
diagnostic calculations, records tool versions and exact commands, and
builds this Markdown into \texttt{P16.5-paper.tex} and
\texttt{P16.5-paper.pdf} with the existing Pandoc and pdfLaTeX. No tools
are installed. The original input and all proof/evidence/log artifacts
remain in this supplied workspace.

The diagnostic program verifies only rational arithmetic in the
displayed examples and the monotonicity inference diagnostic; the proofs
of the geometric and analytic assertions are the arguments above. Source
reading, byte identity, successful program execution and TeX compilation
have no mathematical certification force. Public renderer responses are
saved as such. Their altered HTML rendering dates do not replace the
literal supplied PDF version/date binding. There is no PDF binary-source
checksum unless that binary has actually been acquired; none is asserted
here. No other worker is read. No heavy computation is used, and no
supervised compute facility was exposed or needed.

\end{document}
